\chapter{Quellen}

\sloppy Herangezogene Texte, in der Reihenfolge ihrer Erscheinungsjahre.

\begin{itemize}
\item Baumann, Julius: Die Lehren von Raum, Zeit und Mathematik in der neueren Philosophie nach ihrem ganzen Einfluss dargestellt und beurtheilt. Bd.\,2: Leibniz, Leibniz und Clarke, Berkeley, Hume; Kurzer Lehrbegriff; Schluss und Regeln. Berlin 1869.
\item Schneider, Ilse: Das Raum-Zeit-Problem bei Kant und Einstein. Berlin 1921.
\item Reichenbach, Hans: Philosophie der Raum-Zeit-Lehre (1928). Gesammelte Werke, Bd.\,2, mit Erläuterungen von Andreas Kamlah. Braunschweig 1977.
\item Fink, Eugen: Zur ontologischen Frühgeschichte von Raum -- Zeit -- Bewegung. Den Haag 1957.
\item Brentano, Franz: Philosophische Untersuchungen zu Raum, Zeit und Kontinuum. Aus dem Nachlass herausgegeben von Stephan Körner und Roderick M. Chisholm, mit Anmerkungen von Alfred Kastil. Philosophische Bibliothek Bd.\,293. Hamburg 1976.
\item Bonsiepen, Wolfgang: Hegels Raum-Zeit-Lehre. Dargestellt anhand zweier Vorlesungsnachschriften. Mit Edition der Nachschrift. Hegel-Studien 20, 1985.
\item Koch, Anton Friedrich: Wozu noch Erste Philosophie? Zeitschrift für philosophische Forschung 48, 1994, Heft 4.
\item Koch, Anton Friedrich: Hegelsche Subjekte in Raum und Zeit. Hegel-Jahrbuch 2010.
\item Chiba, Kiyoshi: Kants Ontologie der raumzeitlichen Wirklichkeit. Versuch einer anti-realistischen Interpretation der Kritik der reinen Vernunft. Kantstudien-Ergänzungshefte. Berlin 2012.
\item Guzzoni, Ute: Weile und Weite. Zur nicht-metrischen Erfahrung von Zeit und Raum. Freiburg 2017.
\item Martić, Marko: Zeit und Raum. Eine Untersuchung zu Kants Auffassung der Anschauungsformen. Kantstudien-Ergänzungshefte 217. Berlin 2022.
\end{itemize}

Primärstellen, die über die genannten Texte herangezogen wurden: Aristoteles, Physik IV und De caelo II; Hegel, Enzyklopädie der philosophischen Wissenschaften (1830), \S\S\,254--261, sowie die Vorlesungsnachschriften nach Bonsiepen; Kant, Kritik der reinen Vernunft (B\,275, B\,291, B\,439), Von dem ersten Grunde des Unterschiedes der Gegenden im Raume (1768), Gedanken von der wahren Schätzung der lebendigen Kräfte (1747); Leibniz, Initia rerum mathematicarum metaphysica, nach Baumann; Heidegger, Sein und Zeit \S\,70, Zeit und Sein, Beiträge zur Philosophie, nach Guzzoni.
